\documentclass{article}
\usepackage{amsmath,amsthm}

\newtheorem{lemma}{Lemma}

\begin{document}

\title{Linear Lower Bound on SOS Degree in 1D}
\author{Maximal One}
\date{}
\maketitle

\section*{Setting}

Let
\[
f_k(x) = (1 - x^2)^k
\]
on the compact interval $[-1,1]$.

Assume a Putinar-type representation:
\[
f_k(x) = \sigma_0(x) + \sigma_1(x)(1 - x^2)
\]
where $\sigma_0, \sigma_1$ are sums of squares.

Define $d(k)$ as the minimal degree required in such a representation.

\begin{lemma}
For all $k \ge 1$, one has
\[
d(k) \ge k.
\]
\end{lemma}

\begin{proof}
Observe that $f_k$ vanishes at $x=1$ with multiplicity $k$.
Indeed,
\[
1 - x^2 = (1-x)(1+x)
\]
so near $x=1$, $f_k(x)$ behaves like $(1-x)^k$ up to nonzero scaling.

Evaluating at $x=1$:
\[
f_k(1) = \sigma_0(1).
\]
Since $f_k(1)=0$ and $\sigma_0$ is a sum of squares,
each squared component must vanish at $x=1$.
Hence $\sigma_0$ vanishes with even multiplicity.

If $\sigma_0$ carries the full multiplicity,
its degree must grow at least linearly in $k$.

Otherwise,
$\sigma_1(x)(1-x^2)$ must contribute.
Since $(1-x^2)$ vanishes with multiplicity $1$,
$\sigma_1$ must vanish with multiplicity at least $k-1$.
Again, being a sum of squares,
its vanishing multiplicity is even,
forcing its degree to grow at least linearly in $k$.

Thus $d(k) \ge k$.
\end{proof}

\end{document}
